\section{Introduction}
External-memory agents can turn the end of one episode into persistent state for the next. A completed trajectory is summarized, stored, retrieved, and injected into later context, often without any parameter update. When terminal feedback selects how that summary is written, the evaluator is therefore part of the state-update mechanism rather than only a scorer. An erroneous label can cause the same experience to be committed through a different reflection objective and remain available to future decisions.

Several neighboring literatures establish pieces of this pipeline without separating its stages. JudgeBench and RewardBench document failures of LLM judges and reward models~\citep{tan2025judgebench,lambert2024rewardbench}. ReasoningBank and Reflexion turn success/failure feedback into reusable textual memory~\citep{ouyang2026reasoningbank,shinn2023reflexion}; Mem-$\alpha$ and AttriMem learn memory-construction policies from outcome or process reward~\citep{wang2025memalpha,li2026attrimem}. Memory Reward Inflation and RoMeRL study reward-driven valuation and reuse of memories that already exist~\citep{asadolahi2026inflation,yang2026romerl}, while recent fragility work measures run variance and task-order sensitivity~\citep{ye2026fragility}. We therefore do not claim judge unreliability, feedback-driven memory construction, memory-reward amplification, or self-improvement variance as new. The missing identification object is earlier: under a fixed construction policy and byte-identical experience, can a terminal-label error select a different durable memory state? A second, distinct question is whether that written difference is actually exposed and behaviorally transported by the released retrieval/interface stack.

We first isolate the write boundary. For each source episode, task and trajectory are held byte-for-byte identical while the terminal label counterfactually selects the released success- versus failure-reflection branch. Because those branches necessarily use different instructions, generic prompt sensitivity is the strongest reduction: on eight disjoint trajectories we compare the reward-mode intervention against semantic-preserving rewordings within the same mode whose lexical perturbation is deliberately larger. The original intervention yields four complete pairs, all divergent. We then preregister a breadth replication over 16 additional sources, balanced 8/8 by original benchmark outcome and spanning 11 intent-template families. A 2,200-token parent run is execution-censored in three failure-branch writes; one uniform 4,096-token repair regenerates all 32 units fresh and completes 16/16 pairs. Combined with the original sample, all 20 complete paired sources change both memory content and title sets. Thus the upstream write-time effect is no longer a four-source phenomenon.

We next separate \emph{intervention sensitivity} from \emph{realized transport}. Under a forced memory swap, the original four complete source pairs alter 7/12 aligned next actions. A same-support $4\times4$ terminal replication after an initial non-pass keeps the same support and dual gate, increases only replication depth, and yields mean absolute success-rate difference 0.15625 with permutation $p=0.00074$. This establishes that the paired persistent states \emph{can} alter later outcomes when experimentally injected under matched evidence. However, exact replay of the released ReasoningBank retriever changes the interpretation of that result: with the original four-memory bank, only 8/188 held-out Shopping tasks pass the released all-MiniLM-L6-v2 top-$1$, 0.3 threshold, and none of the four forced-swap future tasks would retrieve those memories.

We therefore expand the source bank before testing ecological transport. With 20 source descriptions, exact released-retriever exposure rises to 125/172 held-out Shopping tasks (72.7\%), spanning 39 intent-template families. Before any terminal outcomes, we freeze the 36 tasks that additionally have released trajectories and deterministic offline evaluators. For each task, the retriever itself selects the source memory; we then swap only that source's success- versus failure-conditioned memory and serialize it with the released ReasoningBank human-memory wrapper. All 288 terminal calls complete. The broad retrieval-matched contrast is small: mean $|\Delta|=0.02083$ with $p=0.4289$, and 34/36 tasks have zero branch contrast. Eighteen of those are joint 0/0 floors and sixteen joint 1/1 ceilings. A matched 144-call no-memory arm gives mean memory-presence effect 0.04514; its omnibus $p=0.00147$ is statistically detectable but misses the unchanged 0.15 practical-effect floor, with omission equidistant from the two reward-conditioned branches on 34/36 tasks.

These results expose a useful boundary rather than a contradiction. Reward-conditioned writing robustly changes persistent state; a forced swap proves that some such state differences can matter; but the released retrieval/interface path does not guarantee a practically large terminal effect. Retrieval exposure, interface, endpoint headroom, and policy/support are separate gates between state corruption and realized behavior. A GLM-5.3 writer independently reinforces this separation: all four paired memories diverge, while its terminal replication misses the practical-effect floor. We also preregister a DeepSeek-V4-Flash downstream-policy transfer on the same 36 tasks, but both the 900-token parent and the single allowed 2,200-token output-cap repair terminate on the first frozen unit with provider length censoring and no assistant text; zero scientific units are produced, so cross-policy transport remains unresolved rather than negative.

\paragraph{Contribution and closest-work boundary.}
The contribution is an identification-and-transport decomposition for reward-conditioned agent memory. First, a byte-identical write intervention plus a stronger reward-preserving wording reduction establishes and broadens the write-time state channel. Second, forced downstream swaps establish conditional behavioral sensitivity without being mislabeled as native retrieval transport. Third, exact released-retriever audits and a 36-task native-wrapper experiment identify a negative transport boundary: broad exposure can coexist with weak terminal branch and memory-presence effects. We make no claim of a new memory architecture, a new reward-robustness method, universally harmful failure memory, writer- or policy-invariant downstream effects, live-browser transport, or a general law of reward-corruption variance.
